\subsection{\texorpdfstring{Complements on spaces \(L^{p}(X,\mu)\)}{Further results on L spaces\^{}\{p\}(X,\textbackslash mu)}}

PROPOSITION 1. --- Let X be a locally compact topological space. Let \(\mu\) be a positive measure on X and \(p \in [1, +\infty[\). Let \((\mathrm{X}_i)_{i \in \mathrm{I}}\) be a locally countable family (INT, IV, p. 190, § 5, no. 9, def. 7) of pairwise disjoint locally compact subsets of X such that the complement of the union of the \(\mathrm{X}_i\) is locally \(\mu\)-negligible. Let \(\varphi_i\) be the characteristic function of \(\mathrm{X}_i\) and \(\mu_i\) be the measure induced by \(\mu\) on \(\mathrm{X}_i\) (INT, IV, p. 186, §5, n° 7, def. 4).

\begin{enumerate}
\def\labelenumi{\alph{enumi})}
\item
  The map \(p_{i}\) : \(f \mapsto f\varphi_{i}\) is a projector of \(\mathcal{L}^{p}(\mathrm{X}, \mu)\) and, by passage to quotients, defines a projector \(\widetilde{p}_{i}\) of \(\mathrm{L}^{p}(\mathrm{X}, \mu)\) . The latter is an orthogonal projector of \(\mathrm{L}^{2}(\mathrm{X}, \mu)\) if p = 2;
\item
  The map \(f \mapsto f|X_i\) induces an isometric isomorphism from the image of \(p_i\) onto the space \(\mathcal{L}^p(X_i, \mu_i)\) and, by passage to quotients, defines an isometric isomorphism from the image of \(\widetilde{p}_i\) onto \(\mathrm{L}^p(X_i, \mu_i)\) , by means of which these two spaces are identified;
\item
  The sum of the spaces \(\mathrm{L}^{p}(\mathrm{X}_{i},\mu_{i})\) is total in the space \(\mathrm{L}^{p}(\mathrm{X},\mu)\) . In the case p = 2, the space \(\mathrm{L}^{2}(\mathrm{X},\mu)\) is the Hilbert sum of the spaces \(\mathrm{L}^{2}(\mathrm{X}_{i},\mu_{i})\) .
\end{enumerate}

The assertion \(a\) ) is elementary. The assertion \(b\) ) follows from the scholium of INT, V, p. 84, § 7, n° 1.

Let \(q \in ]1, +\infty]\) be the conjugate exponent of \(p\) . One identifies the dual of \(L^p(X, \mu)\) with \(L^q(X, \mu)\) (INT, V, p. 61, § 5, n° 8, th. 4). Let \(f\) be a function in \(\mathcal{L}^q(X, \mu)\) whose class \(\widetilde{f}\) in \(L^q(X, \mu)\) satisfies \(\langle \widetilde{f}, \widetilde{p}_i(\varphi) \rangle = 0\) for all \(i \in I\) and every \(\varphi \in L^p(X, \mu)\) . Since the image of \(p_i\) contains \(\mathcal{K}(X_i)\) , it follows that the measure \((f|X_i) \cdot \mu_i\) on \(X_i\) is zero, which means that the restriction of \(f\) to \(X_i\) is locally \(\mu_i\) -negligible on \(X_i\) (INT, V, p. 46, § 5, n° 3, cor. 2). Since the complement in X of the union of the \(X_i\) is locally \(\mu\) -negligible, one concludes that the function \(f\) is locally \(\mu\) -negligible on X (INT, IV, p. 190, § 5, n° 9 and p. 172, n° 2, prop. 5). The class of \(f\) is then zero in \(L^q(X, \mu)\) (using INT, V, p. 8, § 1, n° 3, lemma 1 and the corollary of proposition 9 when \(p \neq 1\) ). By the Hahn--Banach theorem (EVT, II, p. 46, cor. 1), the first part of assertion \(c\) ) follows.

If \(p = 2\) , the image of \(p_i\) is orthogonal to that of \(p_j\) for all \(i \neq j\) since \(X_i\) and \(X_j\) are then disjoint, whence the last assertion.

PROPOSITION 2. --- Let X be a locally compact space that is countable at infinity. Let \(\mu\) be a positive measure on X. For every \(p \in [1, +\infty[\), the space \(\mathrm{L}^{p}(\mathrm{X},\mu)\) is of countable type.

Let \((\mathrm{U}_{n})_{n\in\mathbf{N}}\) be a sequence of relatively compact open sets of X whose union is equal to X and which satisfy \(\overline{U}_{n}\subset U_{n+1}\) for all \(n\in\mathbf{N}\) (TG, I, p. 68, prop. 15). For every \(n\in\mathbf{N}\), the space \(\mathcal{K}(\mathrm{X},\overline{\mathrm{U}}_{n})\) is identified with a closed subspace of the Banach space \(\mathcal{C}(\overline{\mathrm{U}}_{n})\) (INT, III, p. 40, § 1, n° 1); since the latter space is of countable type (TG, X, p. 25, corollary), the same is true of \(\mathcal{K}(\mathrm{X},\overline{\mathrm{U}}_{n})\) (TG, IX, p. 19, cor., (i)). Let \(F_{n}\) be a countable dense subset of \(\mathcal{K}(\mathrm{X},\overline{\mathrm{U}}_{n})\) .

Let \(f \in \mathcal{L}^p(\mathrm{X},\mu)\) and let \(\varepsilon > 0\) . There exists an integer \(n \in \mathbf{N}\) such that \(\int_{\mathrm{X}-\mathrm{U}_n}|f|^p < \varepsilon / 2\) , and there exists \(g \in \mathcal{F}_n\) such that \(\int_{\overline{\mathrm{U}}_n}|f - g|^p < \varepsilon / 2\) . The union of the classes in \(\mathrm{L}^p(\mathrm{X},\mu)\) of the elements of the sets \(\mathcal{F}_n\) is therefore dense in \(\mathrm{L}^p(\mathrm{X},\mu)\) , which concludes the proof (TG, IX, p. 18, prop. 12).

